(1) 对理想基模定义 $z_R=\pi w_0^2/\lambda$，复参数 $q(z)=z+iz_R$。由 $1/q=1/R-i\lambda/(\pi w^2)$ 可解得 $w(z)=w_0\sqrt{1+(z/z_R)^2}$；故束腰处半径严格最小，为 $w_0$，且 $w(-z)=w(z)$。
(2) 复参数的实虚部分别给波前曲率 $R(z)=z[1+(z_R/z)^2]$（$z\ne0$）与 Gouy 相位 $\zeta(z)=\arctan(z/z_R)$。束腰时 $R\to\infty$ 表明波前局部平坦，$\zeta(0)=0$；远场 $|z|\gg z_R$ 时 $R(z)\sim z$，Gouy 相位趋于 $\pm\pi/2$。
(3) 远场 $w(z)\sim w_0|z|/z_R$，故小角半发散角 $\theta=w_0/z_R=\lambda/(\pi w_0)$。方程 $w(z)=2w_0$ 等价于 $1+(z/z_R)^2=4$，因此 $z=\pm\sqrt3z_R$；两侧对称的光束半径与相位曲率具有不同的奇偶性质。
(4) $w_0\to\infty$ 时 $z_R\to\infty$ 且在有限距离上 $w(z)\approx w_0$、$\theta\to0$，对应近似平面波。束腰越小，横向波矢分布越宽，$\theta\propto1/w_0$；若 $\lambda/(\pi w_0)$ 不再远小于一，必须采用非近轴矢量衍射而非上述标量高斯公式。
